% GENERATED FILE. Edit canonical lessons, then run npm run generate:books.
% canonical-source-hash: fnv1a64:8921e8af57c59173
% canonical-lessons: AR-C438-siyaha, AR-C438-jawla, AR-C438-nuzha, AR-C438-tajawwala, AR-C438-rahhala, AR-R438-review-431-438

\chapter{Sightseeing}
\label{ch:ar-sightseeing}
\hlchaptermodality{438}

\begin{chapteropening}
\textbf{By the end of this chapter:} \emph{I can say tourism, a tour, an outing, to wander around and a great traveller.}

Complete the last lesson of chapter 438: I can say tourism, a tour, an outing, to wander around and a great traveller.
\end{chapteropening}

\section[\texorpdfstring{\ar{سياحة} (siyāḥa)}{سياحة (siyāḥa)}]{\ar{سياحة} (siyāḥa) — tourism}

\label{lesson:AR-C438-siyaha}



\begin{quote}
Before the new one: say the Arabic for to guide, then the Arabic for an itinerary.
\end{quote}

\subsection*{You'll want to know: \ar{سياحة}}
\textbf{\ar{سياحة}} — \emph{siyāḥa} — \textquotedblleft{}tourism\textquotedblright{}.

\begin{quote}
\ar{تعمل} \ar{أختي} \ar{في} \ar{السياحة}.
\end{quote}

\begin{quote}
\emph{taʿmalu ukhtī fī as-siyāḥa.}
\end{quote}

\begin{quote}
My sister works in tourism.
\end{quote}

\begin{cousinweb}
\textbf{\ar{سياحة}} (\emph{siyāḥa}) is built on the native Arabic root \emph{s-y-ḥ}, \textquotedblleft{}to flow; to roam the land\textquotedblright{} (Lane). It is the verbal noun of \emph{sāḥa}, \textquotedblleft{}to travel about\textquotedblright{}.
\end{cousinweb}

\begin{grammarlens}[title={What you've built}]
One more word for getting around.
\end{grammarlens}

\subsection*{Your turn}
\begin{itemize}
\raggedright
  \item \emph{Say it:} \emph{siyāḥa}
  \item \emph{Say it:} \emph{siyāḥa}, once more
  \item \emph{Say it:} \emph{taʿmalu ukhtī fī as-siyāḥa.} — My sister works in tourism.
  \item \emph{Recall:} say the Arabic for to guide, then the Arabic for an itinerary, then say \emph{siyāḥa} again
\end{itemize}

\begin{tcolorbox}[breakable,colback=teal!4,colframe=teal!35!black,title={Before you move on}]
What does \ar{سياحة} mean? (\textquotedblleft{}Tourism\textquotedblright{}.) Say it once more.
\end{tcolorbox}
\section[\texorpdfstring{\ar{جولة} (jawla)}{جولة (jawla)}]{\ar{جولة} (jawla) — a tour; a round}

\label{lesson:AR-C438-jawla}



\begin{quote}
Before the new one: say the Arabic for an itinerary, then the Arabic for tourism.
\end{quote}

\subsection*{You'll want to know: \ar{جولة}}
\textbf{\ar{جولة}} — \emph{jawla} — \textquotedblleft{}a tour; a round\textquotedblright{}.

\begin{quote}
\ar{قمنا} \ar{بجولة} \ar{في} \ar{المدينة} \ar{القديمة}.
\end{quote}

\begin{quote}
\emph{qumnā bi-jawla fī al-madīna al-qadīma.}
\end{quote}

\begin{quote}
We went on a tour of the old city.
\end{quote}

\emph{qāma bi-jawla}: to go on a tour.

\begin{cousinweb}
\textbf{\ar{جولة}} (\emph{jawla}) is built on the native Arabic root \emph{j-w-l}, \textquotedblleft{}to go round, to roam\textquotedblright{} (Lane).
\end{cousinweb}

\begin{grammarlens}[title={What you've built}]
One more word for getting around.
\end{grammarlens}

\subsection*{Your turn}
\begin{itemize}
\raggedright
  \item \emph{Say it:} \emph{jawla}
  \item \emph{Say it:} \emph{jawla}, once more
  \item \emph{Say it:} \emph{qumnā bi-jawla fī al-madīna al-qadīma.} — We went on a tour of the old city.
  \item \emph{Recall:} say the Arabic for an itinerary, then the Arabic for tourism, then say \emph{jawla} again
\end{itemize}

\begin{tcolorbox}[breakable,colback=teal!4,colframe=teal!35!black,title={Before you move on}]
What does \ar{جولة} mean? (\textquotedblleft{}A tour; a round\textquotedblright{}.) Say it once more.
\end{tcolorbox}
\section[\texorpdfstring{\ar{نزهة} (nuzha)}{نزهة (nuzha)}]{\ar{نزهة} (nuzha) — an outing; a pleasure trip}

\label{lesson:AR-C438-nuzha}



\begin{quote}
Before the new one: say the Arabic for tourism, then the Arabic for a tour.
\end{quote}

\subsection*{You'll want to know: \ar{نزهة}}
\textbf{\ar{نزهة}} — \emph{nuzha} — \textquotedblleft{}an outing; a pleasure trip\textquotedblright{}.

\begin{quote}
\ar{خرجنا} \ar{في} \ar{نزهة} \ar{إلى} \ar{البحيرة}.
\end{quote}

\begin{quote}
\emph{kharajnā fī nuzha ilā al-buḥayra.}
\end{quote}

\begin{quote}
We went out on an outing to the lake.
\end{quote}

\begin{cousinweb}
\textbf{\ar{نزهة}} (\emph{nuzha}) is usually traced to the native Arabic root \emph{n-z-h}, \textquotedblleft{}to keep far from what is unclean\textquotedblright{}.
\end{cousinweb}

\begin{grammarlens}[title={What you've built}]
One more word for getting around.
\end{grammarlens}

\subsection*{Your turn}
\begin{itemize}
\raggedright
  \item \emph{Say it:} \emph{nuzha}
  \item \emph{Say it:} \emph{nuzha}, once more
  \item \emph{Say it:} \emph{kharajnā fī nuzha ilā al-buḥayra.} — We went out on an outing to the lake.
  \item \emph{Recall:} say the Arabic for tourism, then the Arabic for a tour, then say \emph{nuzha} again
\end{itemize}

\begin{tcolorbox}[breakable,colback=teal!4,colframe=teal!35!black,title={Before you move on}]
What does \ar{نزهة} mean? (\textquotedblleft{}An outing; a pleasure trip\textquotedblright{}.) Say it once more.
\end{tcolorbox}
\section[\texorpdfstring{\ar{تجول} (tajawwala)}{تجول (tajawwala)}]{\ar{تجول} (tajawwala) — to wander around; to tour}

\label{lesson:AR-C438-tajawwala}



\begin{quote}
Before the new one: say the Arabic for a tour, then the Arabic for an outing.
\end{quote}

\subsection*{You'll want to know: \ar{تجول}}
\textbf{\ar{تجول}} — \emph{tajawwala} — \textquotedblleft{}to wander around; to tour\textquotedblright{}.

\begin{quote}
\ar{تجولنا} \ar{في} \ar{السوق} \ar{حتى} \ar{المساء}.
\end{quote}

\begin{quote}
\emph{tajawwalnā fī as-sūq ḥattā al-masāʾ.}
\end{quote}

\begin{quote}
We wandered around the market until the evening.
\end{quote}

Present: \ar{يتجول}, \emph{yatajawwal}.

\begin{cousinweb}
\textbf{\ar{تجول}} (\emph{tajawwala}) is built on the native Arabic root \emph{j-w-l}, \textquotedblleft{}to go round, to roam\textquotedblright{} (Lane). It is Form V.
\end{cousinweb}

\begin{grammarlens}[title={What you've built}]
One more word for getting around.
\end{grammarlens}

\subsection*{Your turn}
\begin{itemize}
\raggedright
  \item \emph{Say it:} \emph{tajawwala}
  \item \emph{Say it:} \emph{tajawwala}, once more
  \item \emph{Say it:} \emph{tajawwalnā fī as-sūq ḥattā al-masāʾ.} — We wandered around the market until the evening.
  \item \emph{Recall:} say the Arabic for a tour, then the Arabic for an outing, then say \emph{tajawwala} again
\end{itemize}

\begin{tcolorbox}[breakable,colback=teal!4,colframe=teal!35!black,title={Before you move on}]
What does \ar{تجول} mean? (\textquotedblleft{}To wander around; to tour\textquotedblright{}.) Say it once more.
\end{tcolorbox}
\section[\texorpdfstring{\ar{رحالة} (raḥḥāla)}{رحالة (raḥḥāla)}]{\ar{رحالة} (raḥḥāla) — a great traveller; an explorer}

\label{lesson:AR-C438-rahhala}



\begin{quote}
Before the new one: say the Arabic for an outing, then the Arabic for to wander around.
\end{quote}

\subsection*{You'll want to know: \ar{رحالة}}
\textbf{\ar{رحالة}} — \emph{raḥḥāla} — \textquotedblleft{}a great traveller; an explorer\textquotedblright{}.

\begin{quote}
\ar{كتب} \ar{الرحالة} \ar{عن} \ar{رحلته} \ar{في} \ar{الصحراء}.
\end{quote}

\begin{quote}
\emph{kataba ar-raḥḥāla ʿan riḥlatihi fī aṣ-ṣaḥrāʾ.}
\end{quote}

\begin{quote}
The explorer wrote about his journey in the desert.
\end{quote}

The ending \emph{-a} here marks a man who does a thing a great deal, not a woman.

\begin{cousinweb}
\textbf{\ar{رحالة}} (\emph{raḥḥāla}) is built on the native Arabic root \emph{r-ḥ-l}, \textquotedblleft{}to set out on a journey\textquotedblright{} (Lane). It shares its root with \emph{riḥla}, \textquotedblleft{}a journey\textquotedblright{}.
\end{cousinweb}

\begin{grammarlens}[title={What you've built}]
One more word for getting around.
\end{grammarlens}

\subsection*{Your turn}
\begin{itemize}
\raggedright
  \item \emph{Say it:} \emph{raḥḥāla}
  \item \emph{Say it:} \emph{raḥḥāla}, once more
  \item \emph{Say it:} \emph{kataba ar-raḥḥāla ʿan riḥlatihi fī aṣ-ṣaḥrāʾ.} — The explorer wrote about his journey in the desert.
  \item \emph{Recall:} say the Arabic for an outing, then the Arabic for to wander around, then say \emph{raḥḥāla} again
\end{itemize}

\begin{tcolorbox}[breakable,colback=teal!4,colframe=teal!35!black,title={Before you move on}]
What does \ar{رحالة} mean? (\textquotedblleft{}A great traveller; an explorer\textquotedblright{}.) Say it once more.
\end{tcolorbox}
\section[(dialogue)]{Words from chapters 431–438 — a review}

\label{lesson:AR-R438-review-431-438}



\begin{quote}
Say the Arabic for a tunnel and a signal.
\end{quote}

\subsection*{Your turn}
Say these aloud:

\begin{itemize}
\raggedright
  \item a tunnel, a signal, traffic, congestion, to turn off
  \item fuel, petrol, an engine, brakes, to supply
  \item a tyre, a wheel, a hole, a pump, to tow
  \item a licence, a bend, to renew, to issue, to overtake
  \item a route, a fare, a subscription, a pedestrian crossing, to turn around
  \item a garage, a kiosk, a ticket window, a porter, to travel about
  \item a compass, speed, a tour guide, to guide, an itinerary
  \item tourism, a tour, an outing, to wander around, a great traveller
\end{itemize}

\begin{tcolorbox}[breakable,colback=teal!4,colframe=teal!35!black,title={Before you move on}]
A tunnel? (\textbf{\ar{نفق}}, \emph{nafaq}.) A signal? (\textbf{\ar{إشارة}}, \emph{ishāra}.) Traffic? (\textbf{\ar{مرور}}, \emph{murūr}.) Congestion? (\textbf{\ar{ازدحام}}, \emph{izdiḥām}.) To turn off? (\textbf{\ar{انعطف}}, \emph{inʿaṭafa}.) Fuel? (\textbf{\ar{وقود}}, \emph{waqūd}.) Petrol? (\textbf{\ar{بنزين}}, \emph{banzīn}.) An engine? (\textbf{\ar{محرك}}, \emph{muḥarrik}.) Brakes? (\textbf{\ar{فرامل}}, \emph{farāmil}.) To supply? (\textbf{\ar{زود}}, \emph{zawwada}.) A tyre? (\textbf{\ar{إطار}}, \emph{iṭār}.) A wheel? (\textbf{\ar{عجلة}}, \emph{ʿajala}.) A hole? (\textbf{\ar{ثقب}}, \emph{thuqb}.) A pump? (\textbf{\ar{مضخة}}, \emph{miḍakhkha}.) To tow? (\textbf{\ar{جر}}, \emph{jarra}.) A licence? (\textbf{\ar{رخصة}}, \emph{rukhṣa}.) A bend? (\textbf{\ar{منعطف}}, \emph{munʿaṭaf}.) To renew? (\textbf{\ar{جدد}}, \emph{jaddada}.) To issue? (\textbf{\ar{أصدر}}, \emph{aṣdara}.) To overtake? (\textbf{\ar{تجاوز}}, \emph{tajāwaza}.) A route? (\textbf{\ar{مسار}}, \emph{masār}.) A fare? (\textbf{\ar{أجرة}}, \emph{ujra}.) A subscription? (\textbf{\ar{اشتراك}}, \emph{ishtirāk}.) A pedestrian crossing? (\textbf{\ar{ممر} \ar{مشاة}}, \emph{mamarr mushāh}.) To turn around? (\textbf{\ar{استدار}}, \emph{istadāra}.) A garage? (\textbf{\ar{مرآب}}, \emph{mirʾāb}.) A kiosk? (\textbf{\ar{كشك}}, \emph{kushk}.) A ticket window? (\textbf{\ar{شباك}}, \emph{shubbāk}.) A porter? (\textbf{\ar{حمال}}, \emph{ḥammāl}.) To travel about? (\textbf{\ar{تنقل}}, \emph{tanaqqala}.) A compass? (\textbf{\ar{بوصلة}}, \emph{būṣala}.) Speed? (\textbf{\ar{سرعة}}, \emph{surʿa}.) A tour guide? (\textbf{\ar{مرشد}}, \emph{murshid}.) To guide? (\textbf{\ar{أرشد}}, \emph{arshada}.) An itinerary? (\textbf{\ar{خط} \ar{سير}}, \emph{khaṭṭ sayr}.) Tourism? (\textbf{\ar{سياحة}}, \emph{siyāḥa}.) A tour? (\textbf{\ar{جولة}}, \emph{jawla}.) An outing? (\textbf{\ar{نزهة}}, \emph{nuzha}.) To wander around? (\textbf{\ar{تجول}}, \emph{tajawwala}.) A great traveller? (\textbf{\ar{رحالة}}, \emph{raḥḥāla}.)
\end{tcolorbox}
